\documentclass[10pt,a4paper]{amsart}
\usepackage[margin=19mm]{geometry}
\usepackage{amsmath,amssymb,mathtools}
\usepackage[scr=rsfs,cal=euler]{mathalfa}
\usepackage{xfrac}
\usepackage[unicode,hidelinks,bookmarks=false]{hyperref}
\newcommand{\ie}{i.e.\ }
\newcommand{\cf}{cf.\ }
\newcommand{\resp}{respectively\ }
\newcommand{\Subsection}{\subsection}
\providecommand{\nobreakdash}{\nobreak\hbox{-}}
\setlength{\emergencystretch}{2em}
\multlinegap0pt
\usepackage[shortlabels]{enumitem}
\usepackage{amssymb}
\usepackage{mathtools}
\newtheorem{definition}{Definition}
\newtheorem{lemma}{Lemma}
\newtheorem{corollary}{Corollary}
  \def\CC{C}%
  \def\HH{H}%
  \def\HHc{H_C}%
  \def\OO{O}%
\newcommand{\RR}{\mathbb R}
\title{The Four-Point Picard Theorem for Quaternionic Slice Regular Functions}
\author{Guangbin Ren, Xin Zhao}
\date{Source: arXiv:2606.08651v2 (CC BY 4.0)}
\begin{document}
\maketitle

\noindent\emph{Source and licence.} The Four-Point Picard Theorem for Quaternionic Slice Regular Functions. Version arXiv:2606.08651v2 (CC BY 4.0), distributed by arXiv under the Creative Commons Attribution 4.0 International licence, http://creativecommons.org/licenses/by/4.0/, as recorded in the arXiv licence metadata for this exact version and shown on the abstract page https://arxiv.org/abs/2606.08651v2. Full licence notice: \texttt{licenses/arxiv-2606.08651-CC-BY-4.0.txt}.\par\smallskip

\noindent\textit{Editorial scope.} Counted unit: the article's corollary cor:finite-order-Q-data (source0.tex lines 807--813). The article's complete original proof of it is delivered with it (source0.tex lines 815--878, one \texttt{proof} environment of 64 lines). No mathematics is written, repaired or re-derived: the statement and the proof are the article's own lines, in the article's own notation, and no retained line is substituted or re-flowed.  Retained uncounted as the source-local context the proof needs: lines 364--375; lines 636--638; lines 694--696; lines 704--706; lines 759--769; lines 787--793.  Nothing was omitted from any retained span, so the package carries no omission record.  No retained line contains a citation, so the wrapper carries no bibliography and no bibliography entry.  No retained line contains a figure environment, an image-inclusion command or drawing code, so no image is delivered and the package declares no asset dependency.  Editorial formatting only, in addition: an AMS \texttt{amsart} preamble that restates the article's own declarations and macros used by the retained lines, namely the environments \texttt{definition}, \texttt{lemma}, \texttt{corollary} on separate counters, together with the standard packages those lines need.  The retained environments print in this document as Definition 1, Lemma 1, Corollary 1 in order of appearance, whereas the article prints the same items with the numbering of its own full document.  The complete native source of the article as distributed in the arXiv e-print for this exact version is bundled unchanged as \texttt{sources/202348/source0.tex}.
\begin{definition}[Four-point omission datum]\label{def:omission-datum}
A four-point omission datum is a pair
\[
(F,(p_0,p_1,p_2,p_3))
\]
where $F:\CC\to\HHc$ is a real-symmetric entire stem function, the points $p_0,p_1,p_2,p_3\in\HH$ are affinely independent and normalized by $p_0=0$, and each function
\[
Q_j=\langle F-p_j,F-p_j\rangle,
\qquad j=0,1,2,3,
\]
is zero-free on $\CC$.  The datum is called nonconstant if $F$ is nonconstant.
\end{definition}

\begin{lemma}[Constant $Q_j$-data]\label{lem:constant-data}
If $Q_0,Q_1,Q_2,Q_3$ are all constant entire functions, then $F=W+\tau n$ is constant.
\end{lemma}

\begin{lemma}[Zero-free finite-order entire functions]\label{lem:finite-order-zero-free}
Let $Q\in\OO(\CC)^*$ be zero-free and of finite order. Then there exists a complex polynomial $P\in\CC[z]$ such that $Q=e^P$. If $Q(x)>0$ for all $x\in\RR$, then $P$ may be chosen with real coefficients.
\end{lemma}

\begin{lemma}[Simple-zero obstruction]\label{lem:simple-zero}
Let $P\in\CC[z]$ be nonconstant and let $C\in\CC^*$. Then neither $e^P-C$ nor $1-Ce^P$ is an entire square.
\end{lemma}

\begin{lemma}[Finite-order rigidity alternatives]\label{lem:rigidity-alternatives}
Assume that the four functions $Q_j$ are zero-free, finite-order, and arise from an omission datum. Write
\[
Q_j=e^{P_j},\qquad P_j\in\RR[z].
\]
Suppose $P_k$ is nonconstant. Along one of the two real directions $x\to+\infty$ or $x\to-\infty$, one of the following alternatives occurs:
\begin{enumerate}[label=(\Alph*)]
\item $P_k(x)\to+\infty$, and then $P_i=P_k$ for every $i\neq k$.
\item $P_k(x)\to-\infty$, and then $P_i\equiv\log d_{ik}$ for every $i\neq k$.
\end{enumerate}
\end{lemma}

\begin{lemma}[Basepoint invariance of the normal square]\label{lem:basepoint-invariance}
Let $p_0,p_1,p_2,p_3$ be affinely independent and let $V_{\mathrm{aff}}$ be the direction space of their real affine span.  Choose a unit vector $n\in V_{\mathrm{aff}}^\perp$ and write, after the normalization $p_0=0$,
\[
F=W+\tau n,\qquad W\in (V_{\mathrm{aff}})_{\CC}.
\]
If $F-p_k=W_k+\tau_k n_k$ is the decomposition obtained after using $p_k$ as basepoint, where $n_k$ is a unit normal to $V_{\mathrm{aff}}$, then $n_k=\pm n$ and $\tau_k^2=\tau^2$.  Consequently the square in the discriminant identity is independent of the chosen basepoint.
\end{lemma}

\begin{corollary}[Finite-order $Q$-data]\label{cor:finite-order-Q-data}
Let $(F,(p_0,p_1,p_2,p_3))$ be a four-point omission datum in the sense of Definition~\ref{def:omission-datum}. Suppose that each
\[
Q_j(z)=\langle F(z)-p_j,F(z)-p_j\rangle,\qquad j=0,1,2,3,
\]
is of finite order. Then $F$ is constant.
\end{corollary}

\begin{proof}
For real $x$, real-symmetry gives $F(x)\in\HH$, and hence
\[
Q_j(x)=\|F(x)-p_j\|^2>0.
\]
By Lemma~\ref{lem:finite-order-zero-free},
\[
Q_j=e^{P_j},\qquad P_j\in\RR[z].
\]
If all $P_j$ are constant, then all $Q_j$ are constant, and Lemma~\ref{lem:constant-data} implies that $F$ is constant.

Assume that some $P_k$ is nonconstant. By Lemma~\ref{lem:rigidity-alternatives}, either Alternative A or Alternative B holds.  The rebasing at $p_k$ used below preserves the normal square by Lemma~\ref{lem:basepoint-invariance}.

In Alternative A, all $Q_j$ are equal to $U=e^{P_k}$. Use $p_k$ as the basepoint. Write
\[
\{i_1,i_2,i_3\}=\{0,1,2,3\}\setminus\{k\},\qquad
r_\alpha=p_{i_\alpha}-p_k\quad(1\le \alpha\le3),
\]
and let $G_k$ be the Gram matrix of $r_1,r_2,r_3$. Then all four $Q$-functions are $U$, so
\[
b_\alpha=\frac{U+\|r_\alpha\|^2-U}{2}=\frac{\|r_\alpha\|^2}{2}.
\]
With $d_k=(\|r_1\|^2,\|r_2\|^2,\|r_3\|^2)^T$, the discriminant identity gives
\[
\tau^2=U-\frac14d_k^TG_k^{-1}d_k.
\]
The constant $C_k=\frac14d_k^TG_k^{-1}d_k$ is positive, since $G_k^{-1}$ is positive definite and $d_k\neq0$. Hence
\[
\tau^2=e^{P_k}-C_k,
\]
contradicting Lemma~\ref{lem:simple-zero}.

In Alternative B, the other three functions are constants:
\[
Q_i\equiv d_{ik}=\|p_i-p_k\|^2\qquad (i\neq k).
\]
Again using $p_k$ as basepoint, put $U=Q_k=e^{P_k}$. Then
\[
b=\frac U2(1,1,1)^T.
\]
Thus
\[
\tau^2=U-\frac{U^2}{4}(1,1,1)G_k^{-1}(1,1,1)^T.
\]
Let
\[
D_k=\frac14(1,1,1)G_k^{-1}(1,1,1)^T>0.
\]
Then
\[
\tau^2=U(1-D_kU)=e^{P_k}(1-D_ke^{P_k}).
\]
Since $e^{P_k}$ has the entire square root $e^{P_k/2}$, the function
\[
S:=\tau e^{-P_k/2}
\]
is entire and satisfies
\[
S^2=1-D_ke^{P_k},
\]
again contradicting Lemma~\ref{lem:simple-zero}.

Both alternatives are impossible. Hence all $P_j$ are constant, and $F$ is constant.
\end{proof}

\end{document}
